
\documentclass[11pt]{article}
\usepackage[margin=1in]{geometry}
\usepackage{amsmath,amssymb,amsthm,mathtools}
\usepackage{microtype}
\usepackage[colorlinks=true,linkcolor=blue,citecolor=blue,urlcolor=blue]{hyperref}

\newtheorem{definition}{Definition}
\newtheorem{lemma}{Lemma}
\newtheorem{theorem}{Theorem}
\newtheorem{corollary}{Corollary}
\newtheorem{remark}{Remark}

\newcommand{\ip}[2]{\left\langle #1,#2\right\rangle}
\newcommand{\norm}[1]{\left\lVert #1\right\rVert}

\title{A Self-Concordant Level-Set Estimate via the $E\|U\|$ Bound}
\author{}
\date{}

\begin{document}
\maketitle

\section{Definitions}

Let $\phi$ be a $C^3$ strictly convex barrier on a convex domain. Write
\[
    g(x):=\nabla \phi(x),
    \qquad
    H(x):=\nabla^2\phi(x).
\]
The Hessian metric defines the local primal and dual norms
\[
    \|u\|_x^2:=u^\top H(x)u,
    \qquad
    \|w\|_{x,*}^2:=w^\top H(x)^{-1}w.
\]
The corresponding Riemannian distance is denoted by $d_R$.

The Jeffreys divergence is
\[
    V(x,y):=\ip{x-y}{g(x)-g(y)}.
\]
We also use the endpoint residual energy
\[
    d(x,y):=\|x-y\|_x^2+\|g(x)-g(y)\|_{x,*}^2,
\]
and the asymmetry residual
\[
    U(x,y):=(x-y)-H(x)^{-1}\bigl(g(x)-g(y)\bigr).
\]

\begin{definition}[Hessian stability]
Let $\rho:[0,\infty)\to[1,\infty)$ satisfy $\rho(0)=1$. We say that the Hessian metric is $\rho$-stable if, for all $x,z$,
\[
    \rho(d_R(x,z))^{-1}H(x)
    \preceq H(z)
    \preceq \rho(d_R(x,z))H(x).
\]
For a standard self-concordant barrier, the corresponding exponential stability modulus is
\[
    \rho(r)=e^{2r}.
\]
\end{definition}

Throughout the self-concordant part below we use the standard convention
\[
    |D^3\phi(x)[a,b,b]|
    \le
    2\|a\|_x\|b\|_x^2.
\]

\section{The basic integral error bound}

Let $x(t)$, $t\in[0,1]$, be a unit-time geodesic for the Hessian metric. Assume it has constant kinetic energy
\[
    E:=\|\dot x(t)\|_{x(t)}^2.
\]
Fix a target point $y$ and define
\[
    v(t):=V(x(t),y),
    \qquad
    v_0:=v(0).
\]

\begin{lemma}[Exact integral error under descent]
Assume that along the geodesic one has the second-variation identity
\[
    v''(t)
    =
    2E-\ip{\ddot x(t)}{U(x(t),y)}_{x(t)}.
\]
If the descent condition
\[
    v'(0)\le -(E+v_0)
\]
holds, then
\[
    v(1)
    \le
    \int_0^1(1-t)
    \left|
    \ip{\ddot x(t)}{U(x(t),y)}_{x(t)}
    \right|dt.
\]
If, in addition, the barrier is standard self-concordant, then
\[
    \|\ddot x(t)\|_{x(t)}\le E,
\]
and hence
\[
    v(1)
    \le
    E\int_0^1(1-t)\|U(x(t),y)\|_{x(t)}dt.
\]
\end{lemma}

\begin{proof}
Taylor's formula in integral form gives the exact identity
\[
    v(1)=v(0)+v'(0)+\int_0^1(1-t)v''(t)dt.
\]
Using the assumed second-variation identity,
\[
\begin{aligned}
    \int_0^1(1-t)v''(t)dt
    &=
    \int_0^1(1-t)\left(2E-\ip{\ddot x(t)}{U(x(t),y)}_{x(t)}\right)dt \\
    &=
    E-
    \int_0^1(1-t)\ip{\ddot x(t)}{U(x(t),y)}_{x(t)}dt,
\end{aligned}
\]
because $2E\int_0^1(1-t)dt=E$. Therefore
\[
    v(1)
    =
    v_0+v'(0)+E
    -
    \int_0^1(1-t)\ip{\ddot x(t)}{U(x(t),y)}_{x(t)}dt.
\]
The descent condition gives $v_0+v'(0)+E\le0$, so
\[
    v(1)
    \le
    -\int_0^1(1-t)\ip{\ddot x(t)}{U(x(t),y)}_{x(t)}dt.
\]
Taking absolute values gives the first claim.

For the second claim, the geodesic equation and standard self-concordance give
\[
    \|\ddot x(t)\|_{x(t)}\le E.
\]
Thus
\[
    \left|\ip{\ddot x(t)}{U(x(t),y)}_{x(t)}\right|
    \le
    E\|U(x(t),y)\|_{x(t)},
\]
and the result follows.
\end{proof}

\section{Self-concordant exponential bounds}

The next two lemmas provide the two estimates used to turn the integral bound into a scalar level-set recursion.

\begin{lemma}[Exponential bound on $U$]
For a standard self-concordant barrier,
\[
    \|U(x,y)\|_x
    \le
    2\left(\cosh d_R(x,y)-1\right).
\]
Consequently, since $d_R(x,y)^2\le V(x,y)$,
\[
    \|U(x,y)\|_x
    \le
    2\left(\cosh\sqrt{V(x,y)}-1\right).
\]
\end{lemma}

\begin{proof}
Let $\gamma:[0,R]\to\operatorname{dom}\phi$ be a unit-speed geodesic from $x$ to $y$, where $R=d_R(x,y)$. Then
\[
    y-x=\int_0^R \gamma'(s)ds,
    \qquad
    g(y)-g(x)=\int_0^R H(\gamma(s))\gamma'(s)ds.
\]
Therefore
\[
    U(x,y)
    =
    \int_0^R
    \left[I-H(x)^{-1}H(\gamma(s))\right]\gamma'(s)ds,
\]
up to an overall sign, which is irrelevant after taking norms.

Let
\[
    A_s:=H(x)^{-1/2}H(\gamma(s))H(x)^{-1/2}.
\]
Standard self-concordance gives exponential Hessian stability,
\[
    e^{-2s}I\preceq A_s\preceq e^{2s}I.
\]
Let $v_s:=H(x)^{1/2}\gamma'(s)$. Since $\gamma$ has unit speed,
\[
    v_s^\top A_s v_s=\|\gamma'(s)\|_{\gamma(s)}^2=1.
\]
A spectral calculation gives
\[
    \|(A_s-I)v_s\|
    \le
    e^s-e^{-s}.
\]
Indeed, after setting $w=A_s^{1/2}v_s$, so that $\|w\|=1$,
\[
    (A_s-I)v_s=(A_s-I)A_s^{-1/2}w,
\]
and the operator norm of $(A_s-I)A_s^{-1/2}$ is
\[
    \max_{\lambda\in[e^{-2s},e^{2s}]}
    \frac{|\lambda-1|}{\sqrt\lambda}
    =e^s-e^{-s}.
\]
Thus
\[
    \|U(x,y)\|_x
    \le
    \int_0^R(e^s-e^{-s})ds
    =2(\cosh R-1).
\]
Finally, $R^2=d_R(x,y)^2\le V(x,y)$, and the second claim follows by monotonicity of $\cosh$ on $[0,\infty)$.
\end{proof}

\begin{lemma}[Exponential bound on $d(x,y)$]
For a standard self-concordant barrier,
\[
    d(x,y)
    \le
    2\left(e^{\sqrt{V(x,y)}}-1\right)^2.
\]
\end{lemma}

\begin{proof}
Let $R=d_R(x,y)$, and let $\gamma:[0,R]\to\operatorname{dom}\phi$ be a unit-speed geodesic from $x$ to $y$. As above,
\[
    y-x=\int_0^R \gamma'(s)ds,
    \qquad
    g(y)-g(x)=\int_0^R H(\gamma(s))\gamma'(s)ds.
\]
By exponential Hessian stability,
\[
    H(\gamma(s))\succeq e^{-2s}H(x),
\]
so
\[
    \|\gamma'(s)\|_x
    \le
    e^s\|\gamma'(s)\|_{\gamma(s)}=e^s.
\]
Therefore
\[
    \|x-y\|_x
    \le
    \int_0^R e^sds
    =e^R-1.
\]
Similarly,
\[
    \|H(\gamma(s))\gamma'(s)\|_{x,*}
    \le
    e^s\|\gamma'(s)\|_{\gamma(s)}=e^s,
\]
and hence
\[
    \|g(x)-g(y)\|_{x,*}
    \le
    e^R-1.
\]
Thus
\[
    d(x,y)
    =
    \|x-y\|_x^2+
    \|g(x)-g(y)\|_{x,*}^2
    \le
    2(e^R-1)^2.
\]
Using $R^2\le V(x,y)$ gives
\[
    d(x,y)
    \le
    2\left(e^{\sqrt{V(x,y)}}-1\right)^2.
\]
\end{proof}

\section{Main self-concordant consequence}


\begin{lemma}[Sublevel invariance along the geodesic step]
Let \(x(t)\), \(t\in[0,1]\), be a unit-time Hessian geodesic with constant
kinetic energy
\[
    E:=\|\dot x(t)\|_{x(t)}^2 .
\]
Fix \(y\), and define
\[
    v(t):=V(x(t),y),
    \qquad
    v_0:=v(0).
\]
Assume the descent condition
\[
    v'(0)\le -(E+v_0),
\]
and the energy compatibility condition
\[
    E\le d(x(0),y),
\]
where
\[
    d(x,y):=\|x-y\|_x^2+\|g(x)-g(y)\|_{x,*}^2.
\]
Assume \(\phi\) is standard self-concordant. Define
\[
    F(v)
    :=
    2(e^{\sqrt v}-1)^2(\cosh\sqrt v-1)
    =
    \frac{(e^{\sqrt v}-1)^4}{e^{\sqrt v}}.
\]
If
\[
    F(v_0)<v_0,
\]
then
\[
    v(t)<v_0
    \qquad
    \forall t\in(0,1].
\]
In particular, the entire geodesic segment remains in the initial strict
sublevel set of \(V(\cdot,y)\).
\end{lemma}

\begin{proof}
For any \(t\in[0,1]\), Taylor's integral formula gives
\[
    v(t)
    =
    v_0+t v'(0)+\int_0^t(t-s)v''(s)\,ds.
\]
Using
\[
    v''(s)
    =
    2E-\langle \ddot x(s),U(x(s),y)\rangle_{x(s)},
\]
we obtain
\[
\begin{aligned}
    v(t)
    &=
    v_0+t v'(0)+E t^2
    -
    \int_0^t(t-s)
    \langle \ddot x(s),U(x(s),y)\rangle_{x(s)}\,ds.
\end{aligned}
\]
By the descent condition,
\[
    t v'(0)\le -t(E+v_0),
\]
and hence
\[
\begin{aligned}
    v(t)
    &\le
    v_0-tv_0-Et(1-t)
    +
    \int_0^t(t-s)
    \left|
    \langle \ddot x(s),U(x(s),y)\rangle_{x(s)}
    \right|ds.
\end{aligned}
\]

We now estimate the final integral on any interval on which
\(v(s)\le v_0\).  Standard self-concordance and the geodesic equation imply
\[
    \|\ddot x(s)\|_{x(s)}\le E.
\]
Moreover, the exponential residual bound gives
\[
    \|U(z,y)\|_z
    \le
    2\left(\cosh\sqrt{V(z,y)}-1\right).
\]
Thus, whenever \(v(s)\le v_0\),
\[
    \|U(x(s),y)\|_{x(s)}
    \le
    2\left(\cosh\sqrt{v_0}-1\right).
\]
Therefore
\[
\begin{aligned}
    \int_0^t(t-s)
    \left|
    \langle \ddot x(s),U(x(s),y)\rangle_{x(s)}
    \right|ds
    &\le
    E\cdot 2\left(\cosh\sqrt{v_0}-1\right)
    \int_0^t(t-s)\,ds  \\
    &=
    E\left(\cosh\sqrt{v_0}-1\right)t^2.
\end{aligned}
\]
Consequently, as long as \(v(s)\le v_0\) on \([0,t]\),
\[
    v(t)
    \le
    v_0-tv_0-Et(1-t)
    +
    E\left(\cosh\sqrt{v_0}-1\right)t^2.
\]
Dropping the nonpositive term \(-Et(1-t)\) and using \(t^2\le t\), we get
\[
    v(t)
    \le
    v_0
    -
    t v_0
    +
    tE\left(\cosh\sqrt{v_0}-1\right).
\]
By the energy compatibility condition and the self-concordant endpoint bound,
\[
    E\le d(x(0),y)
    \le
    2(e^{\sqrt{v_0}}-1)^2.
\]
Hence
\[
    E\left(\cosh\sqrt{v_0}-1\right)
    \le
    2(e^{\sqrt{v_0}}-1)^2(\cosh\sqrt{v_0}-1)
    =
    F(v_0).
\]
Therefore, as long as \(v(s)\le v_0\) on \([0,t]\),
\[
    v(t)
    \le
    v_0-t\bigl(v_0-F(v_0)\bigr).
\]

Now suppose, for contradiction, that the geodesic segment exits the initial
sublevel set. Since \(v'(0)\le -(E+v_0)<0\) whenever \(v_0>0\), the path
initially enters the strict sublevel set. Let \(\tau>0\) be the first exit
time, so that
\[
    v(\tau)=v_0,
    \qquad
    v(t)<v_0 \quad \text{for }0<t<\tau.
\]
The preceding estimate is valid on \([0,\tau]\), and gives
\[
    v(\tau)
    \le
    v_0-\tau\bigl(v_0-F(v_0)\bigr).
\]
If \(F(v_0)<v_0\), then the right-hand side is strictly less than \(v_0\),
contradicting \(v(\tau)=v_0\). Therefore no such first exit time exists, and
\[
    v(t)<v_0
    \qquad
    \forall t\in(0,1].
\]
\end{proof}



\begin{corollary}[Self-concordant level-set recursion]
Let $\phi$ be a standard self-concordant barrier. Let $x(t)$ be a unit-time geodesic with energy $E$, and let
\[
    v(t):=V(x(t),y),
    \qquad
    v_0:=v(0).
\]
Assume
\[
    v'(0)\le -(E+v_0),
\]
\[
    E\le d(x(0),y),
\]
and
\[
    v(t)
    \le
    v_0
    \qquad
    \forall t\in[0,1].
\]
Then
\[
    v(1)
    \le
    F(v_0),
\]
where
\[
    F(v):=\frac{(e^{\sqrt v}-1)^4}{e^{\sqrt v}}.
\]
Consequently,
\[
    \lim_{v\downarrow0}\frac{F(v)}{v^2}=1.
\]
Thus the estimate is quadratically convergent in the sense that
\[
    v(1)\le v_0^2(1+o(1)).
\]
\end{corollary}

\begin{proof}
By the integral error lemma and the self-concordant acceleration bound,
\[
    v(1)
    \le
    E\int_0^1(1-t)\|U(x(t),y)\|_{x(t)}dt.
\]
Since $v(t)\le v_0$, the exponential bound on $U$ gives
\[
    \|U(x(t),y)\|_{x(t)}
    \le
    2(\cosh\sqrt{v_0}-1).
\]
Therefore
\[
    v(1)
    \le
    E(\cosh\sqrt{v_0}-1).
\]
The assumption $E\le d(x(0),y)$ and the exponential bound on $d$ give
\[
    E
    \le
    2(e^{\sqrt{v_0}}-1)^2.
\]
Hence
\[
    v(1)
    \le
    2(e^{\sqrt{v_0}}-1)^2(\cosh\sqrt{v_0}-1).
\]
Using
\[
    2(\cosh s-1)=\frac{(e^s-1)^2}{e^s}
\]
with $s=\sqrt{v_0}$ gives
\[
    v(1)
    \le
    \frac{(e^{\sqrt{v_0}}-1)^4}{e^{\sqrt{v_0}}}.
\]
Finally,
\[
    \frac{F(v)}{v^2}
    =
    \left(\frac{e^{\sqrt v}-1}{\sqrt v}\right)^4e^{-\sqrt v}
    \to1
    \qquad(v\downarrow0).
\]
\end{proof}

\begin{remark}[Explicit finite constants]
For any $\delta>0$, define
\[
    C_\delta
    :=
    \sup_{0<v\le\delta}\frac{F(v)}{v^2}
    =
    \left(\frac{e^{\sqrt\delta}-1}{\sqrt\delta}\right)^4e^{-\sqrt\delta}.
\]
Then $v_0\le\delta$ implies
\[
    v(1)\le C_\delta v_0^2.
\]
For example,
\[
    C_{1/4}
    =
    \left(\frac{e^{1/2}-1}{1/2}\right)^4e^{-1/2}
    \approx1.730.
\]
Moreover, $F(v)<v$ for
\[
    0<v<0.4671709662\ldots,
\]
so the bound certifies strict decrease on this level range.
\end{remark}

\section{Follow-up: bounds for a general stability modulus}

The same proof gives a general-modulus version once the exponential bounds are replaced by the corresponding modulus integrals.

\begin{definition}[Modulus primitives]
Given a nondecreasing Hessian-stability modulus $\rho$, define
\[
    K_\rho(r):=\sqrt{\rho(r)},
\]
\[
    S_\rho(R):=\int_0^R K_\rho(s)ds,
\]
and
\[
    J_\rho(R):=\int_0^R\left(K_\rho(s)-K_\rho(s)^{-1}\right)ds.
\]
Also define the infinitesimal logarithmic slope
\[
    L_\rho:=\limsup_{r\downarrow0}\frac{\log\rho(r)}{r}.
\]
\end{definition}

\begin{lemma}[General-modulus bounds on $U$ and $d$]
Assume the Hessian metric is $\rho$-stable. Then
\[
    \|U(x,y)\|_x
    \le
    J_\rho(d_R(x,y))
    \le
    J_\rho(\sqrt{V(x,y)}),
\]
and
\[
    d(x,y)
    \le
    2S_\rho(d_R(x,y))^2
    \le
    2S_\rho(\sqrt{V(x,y)})^2.
\]
\end{lemma}

\begin{proof}
The proof is identical to the self-concordant geodesic proof, replacing the exponential bounds $e^{\pm 2s}$ by $\rho(s)^{\pm1}$. Along a unit-speed geodesic $\gamma$ from $x$ to $y$, Hessian stability gives
\[
    \rho(s)^{-1}H(x)
    \preceq
    H(\gamma(s))
    \preceq
    \rho(s)H(x).
\]
The bound on $d(x,y)$ follows from
\[
    \|\gamma'(s)\|_x\le K_\rho(s)
\]
and
\[
    \|H(\gamma(s))\gamma'(s)\|_{x,*}\le K_\rho(s),
\]
which integrate to
\[
    \|x-y\|_x\le S_\rho(d_R(x,y)),
    \qquad
    \|g(x)-g(y)\|_{x,*}\le S_\rho(d_R(x,y)).
\]
For $U$, the same spectral calculation as in the self-concordant proof gives the integrand bound
\[
    K_\rho(s)-K_\rho(s)^{-1},
\]
and hence
\[
    \|U(x,y)\|_x\le J_\rho(d_R(x,y)).
\]
Finally, $d_R(x,y)^2\le V(x,y)$ and monotonicity of $S_\rho,J_\rho$ give the $V$-based versions.
\end{proof}

\begin{corollary}[General-modulus level-set recursion]
Assume the Hessian metric is $\rho$-stable and $L_\rho<\infty$. Let $x(t)$ be a unit-time geodesic with energy $E$, and let
\[
    v(t):=V(x(t),y),
    \qquad
    v_0:=v(0).
\]
Assume
\[
    v'(0)\le -(E+v_0),
\]
\[
    E\le d(x(0),y),
\]
and
\[
    v(t)
    \le
    v_0
    \qquad
    \forall t\in[0,1].
\]
Then
\[
    v(1)
    \le
    F_\rho(v_0),
\]
where
\[
    F_\rho(v)
    :=
    \frac{L_\rho}{2}
    S_\rho(\sqrt v)^2
    J_\rho(\sqrt v).
\]
\end{corollary}

\begin{proof}
The integral error lemma gives
\[
    v(1)
    \le
    \int_0^1(1-t)
    \left|\ip{\ddot x(t)}{U(x(t),y)}_{x(t)}\right|dt.
\]
The finite value of $L_\rho$ gives the acceleration bound
\[
    \|\ddot x(t)\|_{x(t)}
    \le
    \frac{L_\rho}{2}E.
\]
Since $v(t)\le v_0$, the general-modulus bound on $U$ gives
\[
    \|U(x(t),y)\|_{x(t)}
    \le
    J_\rho(\sqrt{v_0}).
\]
Therefore
\[
    v(1)
    \le
    \frac{L_\rho}{2}EJ_\rho(\sqrt{v_0})
    \int_0^1(1-t)dt
    =
    \frac{L_\rho}{4}EJ_\rho(\sqrt{v_0}).
\]
Using $E\le d(x(0),y)$ and the general-modulus bound on $d$,
\[
    E
    \le
    2S_\rho(\sqrt{v_0})^2.
\]
Thus
\[
    v(1)
    \le
    \frac{L_\rho}{2}
    S_\rho(\sqrt{v_0})^2
    J_\rho(\sqrt{v_0}).
\]
\end{proof}

\begin{remark}[Local quadratic rate]
If
\[
    \rho(r)=1+L_\rho r+O(r^2)
    \qquad(r\downarrow0),
\]
then
\[
    S_\rho(\sqrt v)=\sqrt v+O(v),
    \qquad
    J_\rho(\sqrt v)=\frac{L_\rho}{2}v+O(v^{3/2}).
\]
Therefore
\[
    F_\rho(v)=\frac{L_\rho^2}{4}v^2+O(v^{5/2}).
\]
So the general-modulus recursion is quadratically convergent in the usual sense, with asymptotic constant $L_\rho^2/4$.
\end{remark}

\begin{remark}[Exponential modulus]
For $\rho(r)=e^{Lr}$,
\[
    S_\rho(R)=\frac{2}{L}\left(e^{LR/2}-1\right),
\]
and
\[
    J_\rho(R)=\frac{4}{L}\left(\cosh\left(\frac{LR}{2}\right)-1\right).
\]
Thus
\[
    F_L(v)
    =
    \frac{4}{L^2}
    \frac{\left(e^{L\sqrt v/2}-1\right)^4}{e^{L\sqrt v/2}}.
\]
For $L=2$, this recovers
\[
    F_2(v)=\frac{(e^{\sqrt v}-1)^4}{e^{\sqrt v}}.
\]
\end{remark}



\section{Bounds from Hessian stability}

\begin{lemma}[Exponential upper bound on \(V\) by geodesic distance]
Assume the Hessian metric satisfies exponential stability with parameter \(L\):
\[
    e^{-L d_R(x,z)}H(x)
    \preceq
    H(z)
    \preceq
    e^{L d_R(x,z)}H(x)
    \qquad
    \forall x,z.
\]
Then, for all \(x,y\),
\[
    V(x,y)
    \le
    \Psi_L(d_R(x,y)),
\]
where
\[
    \Psi_L(R)
    :=
    2\int_0^R (R-r)e^{Lr/2}\,dr.
\]
Equivalently, for \(L>0\),
\[
    \Psi_L(R)
    =
    \frac{8}{L^2}\left(e^{LR/2}-1\right)
    -
    \frac{4}{L}R.
\]
Thus
\[
    \boxed{
    V(x,y)
    \le
    \frac{8}{L^2}\left(e^{\frac{L}{2}d_R(x,y)}-1\right)
    -
    \frac{4}{L}d_R(x,y).
    }
\]
For standard self-concordance, \(L=2\), this becomes
\[
    \boxed{
    V(x,y)
    \le
    2\left(e^{d_R(x,y)}-1-d_R(x,y)\right).
    }
\]
\end{lemma}

\begin{proof}
Let \(R:=d_R(x,y)\), and let \(\gamma:[0,R]\to\operatorname{dom}\phi\) be a unit-speed geodesic from \(x\) to \(y\), so
\[
    \gamma(0)=x,
    \qquad
    \gamma(R)=y,
    \qquad
    \|\gamma'(s)\|_{\gamma(s)}=1.
\]
We have
\[
    y-x=\int_0^R \gamma'(a)\,da
\]
and
\[
    g(y)-g(x)
    =
    \int_0^R H(\gamma(b))\gamma'(b)\,db.
\]
Therefore
\[
\begin{aligned}
    V(x,y)
    &=
    \langle y-x,g(y)-g(x)\rangle \\
    &=
    \int_0^R\int_0^R
    \left\langle
        \gamma'(a),
        H(\gamma(b))\gamma'(b)
    \right\rangle
    da\,db .
\end{aligned}
\]
By Cauchy--Schwarz in the metric \(H(\gamma(b))\),
\[
    \left\langle
        \gamma'(a),
        H(\gamma(b))\gamma'(b)
    \right\rangle
    \le
    \|\gamma'(a)\|_{\gamma(b)}
    \|\gamma'(b)\|_{\gamma(b)}.
\]
Since \(\gamma\) is unit speed,
\[
    \|\gamma'(b)\|_{\gamma(b)}=1.
\]
Exponential stability gives
\[
    H(\gamma(b))
    \preceq
    e^{L d_R(\gamma(a),\gamma(b))}H(\gamma(a)).
\]
Because \(\gamma\) is unit speed,
\[
    d_R(\gamma(a),\gamma(b))\le |a-b|.
\]
Thus
\[
    \|\gamma'(a)\|_{\gamma(b)}
    \le
    e^{L|a-b|/2}
    \|\gamma'(a)\|_{\gamma(a)}
    =
    e^{L|a-b|/2}.
\]
Hence
\[
    V(x,y)
    \le
    \int_0^R\int_0^R e^{L|a-b|/2}\,da\,db.
\]
Using the standard identity
\[
    \int_0^R\int_0^R f(|a-b|)\,da\,db
    =
    2\int_0^R (R-r)f(r)\,dr,
\]
with \(f(r)=e^{Lr/2}\), we obtain
\[
    V(x,y)
    \le
    2\int_0^R (R-r)e^{Lr/2}\,dr
    =
    \Psi_L(R).
\]
Finally, evaluating the integral gives
\[
    \Psi_L(R)
    =
    \frac{8}{L^2}\left(e^{LR/2}-1\right)
    -
    \frac{4}{L}R.
\]
\end{proof}

\end{document}
